% !TEX root = ./main.tex

\ustcsetup{
  title               = {在ATLAS实验上\\基于双轻子和不可见末态\\精确测量ZZ玻色子对产生截面\\并探测暗物质},
  title*              = {Measurement of ZZ Production Cross-section and Searching for Dark Matter 
  %at $\sqrt{S}$$ = 13$~TeV 
  in ATLAS with $\ell^+\ell^-+E^{miss}_{T}$ Final State},
  author              = {魏传顺},
  student-id          = {BA19004051},
  discipline          = {粒子物理与原子核物理},
  supervisor          = {赵政国~教授, 周冰~教授},
  % date                = {2026-06-09},  % 提交日期，默认为今日
  %
  % 研究生需要填写
  % practice-supervisor = {洪专~教授, 李实~教授},  % 工程/专业类学位的实践导师
  clc                 = {O572.21},  % 中图分类号，参考 https://www.clcindex.com/category/
  % 当前clc分类为： O 数理科学和化学  O4 物理学  O57 原子核物理学、高能物理学  O572 高能物理学  O572.2 粒子物理学  O572.21 实验与测定
  udc                 = {539.1},  % UDC 分类号，参考 https://udcsummary.info/php/index.php?lang=chi
  % 当前udc分类为： 539.1 核物理、原子物理学、分子物理
  department          = {近代物理系},
  research-direction  = {高能物理实验},
  %
  % secret-level        = {秘密},     % 绝密|机密|秘密|控阅，注释本行则公开
  % secret-year         = {10},      % 保密/控阅期限
  %
  % 数学字体
  % math-style          = GB,  % 可选：GB, TeX, ISO
  math-font           = xits,  % 可选：stix, xits, libertinus
}


% 加载宏包

% 定理类环境宏包
\usepackage{amsthm}

% 插图
\usepackage{graphicx}

% 英文图题、表题
\usepackage{bicaption}

% 三线表
\usepackage{booktabs}

% 表注
\usepackage{threeparttable}

% 跨页表格
\usepackage{longtable}

% 算法
\usepackage{algorithm}
\usepackage{algorithmic}

% SI 量和单位
\usepackage{siunitx}

% 参考文献使用 BibTeX + natbib 宏包
% 顺序编码制
% \usepackage[sort]{natbib}
% \bibliographystyle{ustcthesis-numeric}

% 著者-出版年制
% \usepackage{natbib}
% \bibliographystyle{ustcthesis-authoryear}

% 参考文献使用 BibLaTeX 宏包
\usepackage[style=ustcthesis-numeric]{biblatex}
% \usepackage[style=ustcthesis-numeric, inline=true]{biblatex}
% \usepackage[style=ustcthesis-authoryear]{biblatex}

% 声明 .bib 数据库
\addbibresource{bib/ustc.bib}

% 配置图片的默认目录
\graphicspath{{figures/}}

% 数学命令
\ExplSyntaxOn
\NewDocumentCommand \dif {}  % 微分符号
  {
    \mathop { } \!
    \str_if_eq:VnTF \l__ustc_math_style_str { TeX }
      { d }
      { \mathrm { d } }
  }
% Euler's number or upright e
\NewDocumentCommand \eu {}
  {
    \str_if_eq:VnTF \l__ustc_math_style_str { TeX }
      { e }
      { \symup { e } }
  }
% Imaginary unit or upright i
\NewDocumentCommand \iu {}
  {
    \str_if_eq:VnTF \l__ustc_math_style_str { TeX }
      { i }
      { \symup { i } }
  }
\ExplSyntaxOff

% 用于写文档的命令
\DeclareRobustCommand\cs[1]{\texttt{\char`\\#1}}
\DeclareRobustCommand\env[1]{\texttt{#1}}
\DeclareRobustCommand\pkg[1]{\textsf{#1}}
\DeclareRobustCommand\file[1]{\nolinkurl{#1}}

% hyperref 宏包在最后调用
\usepackage{hyperref}
